% ═══════════════════════════════════════════════════════════════════
% Paper VII: Universal Hebbian Remodeling —
%   Wolff, Glagov, and Hebb as Special Cases of C2
% ═══════════════════════════════════════════════════════════════════
\documentclass[aps,pre,twocolumn,superscriptaddress,floatfix]{revtex4-2}

\usepackage{amsmath,amssymb,amsthm}
\usepackage{graphicx}
\usepackage{hyperref}
\usepackage{bm}
\usepackage{booktabs}
\usepackage{siunitx}
\usepackage{multirow}

\newtheorem{theorem}{Theorem}
\newtheorem{corollary}{Corollary}
\newtheorem{definition}{Definition}
\newtheorem{remark}{Remark}
\newtheorem{proposition}{Proposition}

\begin{document}

\title{Universal Hebbian Remodeling:\\
Wolff's Law, Glagov's Remodeling, Immune Memory,\\
and Skin Keratinization as Special Cases of\\
$\Delta Z = -\eta\,\Gamma\,x_{\text{pre}}\,x_{\text{post}}$}

\author{Hsi-Yu Huang}
\email{llc.y.huangll@gmail.com}
\affiliation{Independent Researcher, Taiwan}

\date{\today}

% ============================================================
\begin{abstract}
Paper~VI established that the reflection coefficient
$\Gamma = (Z_2 - Z_1)/(Z_2 + Z_1)$ is a thermodynamic
necessity at every impedance interface.
Here we show that constraint~C2 of the $\Gamma$-Net
framework---the Hebbian impedance update
$\Delta Z = -\eta\,\Gamma\,x_{\text{pre}}\,x_{\text{post}}$---is
likewise universal: it is the unique gradient-descent rule
that minimizes reflected energy at a repeatedly stimulated
interface.
We identify nineteen $\Gamma$-networks in the
human body---seventeen distributed tissue networks
(neural, vascular, skeletal, muscular, immune, epithelial,
pulmonary, gastrointestinal, endocrine, lymphatic,
renal, hepatic, connective, adipose, hematopoietic,
cartilage, reproductive)
and two central controllers (cardiac pump, regulatory
brain)---and prove that the established remodeling
law of each is algebraically identical to~C2 under
the appropriate substitution of effort, flow, and impedance
variables.
We further show that cartilage ($\eta \approx 0$) serves
as a falsifiable null case: the predicted irreversible
degradation under chronic $\Gamma > 0$ matches
osteoarthritis exactly.
Concretely: \textbf{Hebb's Law} (synaptic potentiation)
is C2 with $Z = Z_{\text{syn}}$;
\textbf{Wolff's Law} (bone remodeling along stress lines)
is C2 with $Z = Z_{\text{mech}}$;
\textbf{Glagov's remodeling} (arterial wall thickening under
shear stress) is C2 with $Z = Z_{\text{vasc}}$;
\textbf{immune memory} (affinity maturation and memory-cell
formation) is C2 with $Z = Z_{\text{receptor}}$;
\textbf{skin keratinization} (callus formation) is C2
with $Z = Z_{\text{mech,skin}}$;
\textbf{Laplace--Starling remodeling} (cardiac hypertrophy)
is C2 with $Z = Z_{\text{heart}}$;
\textbf{baroreceptor resetting} (allostatic adaptation)
is C2 with $Z = Z_{\text{setpoint}}$;
\textbf{tubuloglomerular feedback} (renal autoregulation)
is C2 with $Z = Z_{\text{renal}}$;
\textbf{liver regeneration} is C2 with $Z = Z_{\text{hepatic}}$;
and \textbf{cartilage} demonstrates the predicted C2 null
case ($\eta \approx 0 \Rightarrow$ irreversible degeneration).

We validate the network-propagation prediction using NHANES
10-cycle data ($n = 49{,}774$): adding blood-pressure
measurements (vascular~$Z$ input) with zero fitted parameters
improves all-cause mortality prediction from
AUC~$= 0.629$ to $0.705$ ($\Delta = +0.075$),
and---critically---improves prediction of \emph{cardiac}
death ($\Delta\text{AUC} = +0.166$), \emph{neurological}
death ($+0.061$), and \emph{renal} death ($+0.024$) through
pure network propagation, not direct measurement.
This multi-organ improvement from a single vascular input
is the empirical signature of inter-network $\Gamma$-coupling
that no single-organ textbook formula can produce.
Together with Paper~VI, which established the universality
of~C1, this result shows that both dynamical constraints
of $\Gamma$-Net emerge from the same variational principle
rather than from ad~hoc modelling choices.

\medskip\noindent
\textbf{Key result.}
Evolution did not invent Hebb's Law for neurons and a
different law for bones.
It discovered the same law once---$\Delta Z = -\eta\,\Gamma\,
x_{\text{pre}}\,x_{\text{post}}$---and expressed it in
nineteen distinct biological contexts: seventeen distributed
tissue networks and two central controllers.
The difference between tissues is not the algorithm but
the \emph{material substrate} through which the algorithm
is physically realized.

\medskip\noindent
\textbf{Circularity disclosure.}
The C2 universality claim is structural (algebraic
equivalence of known remodeling laws).  The NHANES
validation uses external epidemiological data with
zero fitted parameters.
\end{abstract}

\maketitle


% ============================================================
%  SECTION 1 — INTRODUCTION
% ============================================================
\section{Introduction}
\label{sec:intro}

Every living tissue remodels.
Bones thicken along load-bearing axes~\cite{Wolff1892}.
Arteries thicken under sustained hypertension~\cite{Glagov1987}.
Synapses strengthen when pre- and post-synaptic neurons
fire together~\cite{Hebb1949}.
Skin forms calluses where friction is chronic~\cite{Thomas2010}.
The immune system forms memory cells after antigen
exposure~\cite{Ahmed1996}.
Muscles hypertrophy under repeated loading~\cite{Goldberg1975}.

These phenomena are studied by six different disciplines
(neuroscience, vascular biology, orthopaedics, dermatology,
immunology, exercise physiology) with six different
mathematical models.
Yet the qualitative pattern is identical in every case:
\emph{repeated stimulation at an interface drives the
interface toward impedance matching.}

Papers~I--VI~\cite{PaperI,PaperII,PaperIII,PaperIV,PaperV,PaperVI}
of this series derived all $\Gamma$-Net behaviour from
three constraints:
\begin{itemize}
  \item[\textbf{C1}] Energy conservation:
    $\Gamma^2 + T = 1$.
  \item[\textbf{C2}] Hebbian update:
    $\Delta Z = -\eta\,\Gamma\,x_{\text{pre}}\,x_{\text{post}}$.
  \item[\textbf{C3}] Signal protocol:
    all inter-module values carry impedance metadata.
\end{itemize}
Paper~VI proved that C1 is a thermodynamic necessity and
validated the framework against NHANES data.
\emph{This paper proves that C2 is equally universal}:
it is the gradient-descent rule that any adaptive system
must obey if it minimizes reflected energy.
In this sense, Papers~VI and~VII complete the thermodynamic
foundation of $\Gamma$-Net: C1 fixes how signals reflect
at interfaces, and C2 fixes how interfaces must adapt
under repeated drive.

The plan is as follows.
Section~\ref{sec:c2-derivation} derives C2 from the MRP
variational principle.
Section~\ref{sec:nineteen-networks} maps nineteen
remodeling laws onto C2, distinguishing seventeen distributed
tissue networks and two central controllers,
including a falsifiable null case (cartilage).
Section~\ref{sec:nhanes-proof} presents the NHANES
network-propagation evidence.
Section~\ref{sec:evolution} interprets tissue remodeling
as evolutionary convergence on C2.
Section~\ref{sec:predictions} lists falsifiable predictions.
Section~\ref{sec:discussion} discusses implications.


% ============================================================
%  SECTION 2 — DERIVATION OF C2 FROM MRP
% ============================================================
\section{C2 as Gradient Descent on Reflected Energy}
\label{sec:c2-derivation}

\subsection{The variational problem}

The Minimum Reflection Principle states:
\begin{equation}
  A[\Gamma] = \int_0^T \sum_i \Gamma_i^2(t)\,dt \to \min\,,
  \label{eq:MRP}
\end{equation}
where $\Gamma_i = (Z_{L,i} - Z_{0,i})/(Z_{L,i} + Z_{0,i})$
is the reflection coefficient at interface~$i$.

The system can adjust $Z_{L,i}$ (Hebbian remodeling) but not
$Z_{0,i}$ (source impedance, set by upstream physics).

\subsection{Gradient computation}

The instantaneous reflected power at interface~$i$ is
$\Gamma_i^2$.  The gradient with respect to
$Z_{L,i}$ is:
\begin{align}
  \frac{\partial \Gamma_i^2}{\partial Z_{L,i}}
  &= 2\Gamma_i \cdot \frac{\partial \Gamma_i}{\partial Z_{L,i}} \notag \\
  &= 2\Gamma_i \cdot \frac{2Z_{0,i}}{(Z_{L,i}+Z_{0,i})^2}\,.
  \label{eq:gradient}
\end{align}
The gradient-descent update is:
\begin{equation}
  \Delta Z_{L,i} = -\frac{\eta}{2}
    \frac{\partial \Gamma_i^2}{\partial Z_{L,i}}
  = -\eta\,\Gamma_i \cdot
    \frac{2Z_{0,i}}{(Z_{L,i}+Z_{0,i})^2}\,.
  \label{eq:gd-exact}
\end{equation}

\subsection{Signal identification}

At the interface, the incident signal has amplitude
$x_{\text{inc}}$ and splits into reflected and
transmitted parts:
\begin{align}
  x_{\text{ref}} &= \Gamma_i \, x_{\text{inc}}\,, \\
  x_{\text{trans}} &= (1-\Gamma_i^2)^{1/2}\, x_{\text{inc}}\,.
\end{align}
The pre-synaptic signal is the incident wave arriving
at the interface:
\begin{equation}
  x_{\text{pre}} \equiv x_{\text{inc}}\,,
\end{equation}
and the post-synaptic signal is the transmitted wave
propagating into the load:
\begin{equation}
  x_{\text{post}} \equiv x_{\text{trans}}\,.
\end{equation}
In the small-mismatch limit ($|\Gamma_i| \ll 1$), the
denominator factor in Eq.~\eqref{eq:gd-exact} approaches
$1/(2Z_0)$, and $x_{\text{post}} \approx x_{\text{pre}}
\cdot \sqrt{T} \approx x_{\text{pre}}$.  The exact
gradient descent reduces to:
\begin{equation}
  \boxed{\Delta Z_i = -\eta\,\Gamma_i\,
    x_{\text{pre},i}\,x_{\text{post},i}}
  \label{eq:C2}
\end{equation}
This is constraint~C2.  It is not a postulate---it is
the \emph{unique} first-order gradient-descent rule for
minimizing reflected energy at a stimulated interface.


% ============================================================
%  SECTION 3 — TWELVE Γ-NETWORKS
% ============================================================
\section{Nineteen $\Gamma$-Networks: Seventeen Tissues and Two Controllers}
\label{sec:nineteen-networks}

We now demonstrate that the known remodeling law of each
tissue type is algebraically equivalent to C2 under the
appropriate substitution of physical variables.
Critically, we distinguish two architectural levels:
\emph{distributed tissue networks} (Sections~\ref{sec:hebb}--\ref{sec:lymphatic}),
which adapt local impedances along network branches; and
\emph{central controllers} (Sections~\ref{sec:cardiac}--\ref{sec:regulatory-brain}),
which set the source impedance $Z_{\text{source}}$ or
regulatory setpoint $Z_{\text{setpoint}}$ for entire
network trees.

Table~\ref{tab:nineteen-networks} summarizes the mapping.
We prove equivalence for all nineteen networks in detail.

\begin{table*}[t]
\centering
\caption{%
Nineteen $\Gamma$-networks: seventeen distributed tissues and two central controllers.
Every row maps to C2: $\Delta Z = -\eta\,\Gamma\,x_{\text{pre}}\,x_{\text{post}}$.
}
\label{tab:nineteen-networks}
\begin{tabular}{lllllll}
\toprule
Network & Flow ($x$) & Effort & $Z$ definition
  & $x_{\text{pre}}$ & $x_{\text{post}}$ & Known name \\
\midrule
Neural
  & Ion current & Membrane voltage & $R_m + 1/j\omega C_m$
  & Pre-syn.\ spike & Post-syn.\ EPSP & \textbf{Hebb's Law} \\
Vascular
  & Blood flow & Blood pressure & $P/Q$
  & Wall shear stress & Endothelial response & \textbf{Glagov remodeling} \\
Skeletal
  & Mechanical force & Stress & Force/velocity
  & Applied load & Bone strain & \textbf{Wolff's Law} \\
Muscular
  & Force & Contractile stress & Stiffness
  & Neural drive & Muscle tension & \textbf{Davis's Law} \\
Immune
  & Cytokines & Antigen signal & Receptor affinity
  & Antigen (APC) & T/B-cell activation & \textbf{Immune memory} \\
Epithelial
  & Heat + pressure & Mechanical stress & Thermal + mech.\ $Z$
  & External friction & Tissue deformation & \textbf{Keratinization} \\
Pulmonary
  & Gas flow & Airway pressure & $P_{\text{aw}}/\dot{V}$
  & Ventilatory demand & Alveolar expansion & \textbf{Altitude adaptation} \\
Gastrointestinal
  & Nutrients & Osmotic gradient & Absorptive $Z$
  & Luminal content & Epithelial uptake & \textbf{Intestinal adaptation} \\
Endocrine
  & Hormones & Receptor occupancy & $1/K_d$
  & Hormone conc.\ & Intracellular resp.\ & \textbf{Receptor regulation} \\
Lymphatic
  & Immune cells & Interstitial pressure & Drainage $Z$
  & Tissue load & Lymph node response & \textbf{Lymph node hypertrophy} \\
Renal
  & Filtrate flow & Glomerular pressure & $P_{\text{glom}}/Q_{\text{filt}}$
  & Macula densa signal & Arteriolar response & \textbf{TG feedback} \\
Hepatic
  & Metabolites & Portal pressure & Clearance $Z$
  & Portal factors (HGF) & Hepatocyte prolif.\ & \textbf{Liver regeneration} \\
Connective
  & Mechanical force & Tensile stress & $E \cdot A/L$
  & Mechanical strain & Fibroblast response & \textbf{Tendon remodeling} \\
Adipose
  & Energy substrates & Insulin signal & Storage $Z$
  & Insulin signal & Adipocyte response & \textbf{Leptin signaling} \\
Hematopoietic
  & O$_2$ delivery & Tissue $P_{O_2}$ & $P_{O_2}/Q_{O_2}$
  & EPO signal & Erythroid response & \textbf{Erythropoiesis} \\
Cartilage
  & Compressive force & Joint stress & $H_A \cdot A/h$
  & Mechanical load & (Negligible: $\eta \approx 0$) & \textbf{C2 null case} \\
Reproductive
  & FSH / LH & Gonadotropin signal & $1/(N_R \cdot K_a)$
  & Gonadotropin & Granulosa / germ cell & \textbf{Follicular selection} \\
\midrule
\multicolumn{7}{c}{\emph{Central controllers}} \\
\midrule
Cardiac (pump)
  & Ejection flow & Ventricular pressure & $P_{\text{vent}}/Q_{\text{ej}}$
  & Wall stress $\sigma$ & Myocardial strain $\varepsilon$ & \textbf{Laplace--Starling} \\
Regulatory brain
  & Efferent drive & Afferent signal & $Z_{\text{setpoint}}$
  & Sensory afferent & Motor efferent & \textbf{Baroreceptor resetting} \\
\bottomrule
\end{tabular}
\end{table*}


% ------ 3.1 NEURAL: HEBB ------
\subsection{Neural: Hebb's Law (1949)}
\label{sec:hebb}

Hebb's postulate~\cite{Hebb1949}:
``When an axon of cell~A is near enough to excite cell~B
and repeatedly or persistently takes part in firing it,
some growth process or metabolic change takes place
in one or both cells such that A's efficiency, as one
of the cells firing B, is increased.''

Modern formulation (rate-based Hebbian LTP):
\begin{equation}
  \Delta w_{AB} = \eta\,r_A\,r_B\,,
  \label{eq:hebb-classical}
\end{equation}
where $r_A$ and $r_B$ are firing rates and $w_{AB}$ is
synaptic weight.

In impedance language, synaptic weight is inversely related
to synaptic impedance: higher $w$ means lower $Z_{\text{syn}}$
(more current passes through).  Therefore:
\begin{equation}
  \Delta Z_{\text{syn}} \propto -\Delta w
  = -\eta\,r_A\,r_B\,.
\end{equation}
Identifying $r_A = x_{\text{pre}}$, $r_B = x_{\text{post}}$,
and noting that the factor driving the change is the
mismatch $\Gamma$ at the synapse (the synapse changes
\emph{because} there is a mismatch to correct):
\begin{equation}
  \Delta Z_{\text{syn}}
  = -\eta\,\Gamma_{\text{syn}}\,
    x_{\text{pre}}\,x_{\text{post}}\,.
  \quad \checkmark
  \label{eq:hebb-C2}
\end{equation}

\begin{remark}
The asymmetric form of STDP (spike-timing-dependent
plasticity) arises from the causal structure of
$x_{\text{pre}}$ preceding $x_{\text{post}}$:
the correlation is positive when the pre-synaptic spike
precedes the post-synaptic spike (LTP; $\Gamma > 0$)
and negative when reversed (LTD; $\Gamma < 0$).
\end{remark}


% ------ 3.2 SKELETAL: WOLFF ------
\subsection{Skeletal: Wolff's Law (1892)}
\label{sec:wolff}

Wolff's Law~\cite{Wolff1892}:
``Every change in the form and function of a bone, or of
its function alone, is followed by certain definite changes
in its internal architecture.''

Modern mechanotransduction formulation~\cite{Frost2003}:
\begin{equation}
  \frac{d(\rho_{\text{bone}})}{dt} =
  \begin{cases}
    +B(\varepsilon - \varepsilon_{\text{ref}})
      & \text{if } \varepsilon > \varepsilon_{\text{ref}} \\
    -R(\varepsilon_{\text{ref}} - \varepsilon)
      & \text{if } \varepsilon < \varepsilon_{\text{ref}}
  \end{cases}
  \label{eq:wolff-frost}
\end{equation}
where $\varepsilon$ is local strain and $\varepsilon_{\text{ref}}$
is the homeostatic setpoint.

In impedance language, bone mechanical impedance
$Z_{\text{bone}} = F/v$ (force per velocity) is proportional
to density and stiffness.  The strain mismatch
$\varepsilon - \varepsilon_{\text{ref}}$ is
the mechanical reflection coefficient:
\begin{equation}
  \Gamma_{\text{mech}} =
  \frac{Z_{\text{bone}} - Z_{\text{load}}}
       {Z_{\text{bone}} + Z_{\text{load}}}\,,
\end{equation}
where $Z_{\text{load}}$ is the mechanical impedance of
the applied force (ground reaction, muscle pull, gravity).

The applied force is $x_{\text{pre}}$ (the stimulus arriving
at the bone interface); the resulting bone strain is
$x_{\text{post}}$ (the signal transmitted through the bone).
Wolff's Law becomes:
\begin{equation}
  \Delta Z_{\text{bone}}
  = -\eta\,\Gamma_{\text{mech}}\,
    F_{\text{applied}}\,\varepsilon_{\text{bone}}
  = -\eta\,\Gamma\,x_{\text{pre}}\,x_{\text{post}}\,.
  \quad \checkmark
  \label{eq:wolff-C2}
\end{equation}

\begin{remark}
Astronaut bone loss in microgravity:
$F_{\text{applied}} \approx 0$ $\Rightarrow$
$x_{\text{pre}} \cdot x_{\text{post}} \approx 0$ $\Rightarrow$
$\Delta Z = 0$.
Without the Hebbian drive, bone impedance drifts freely
under metabolic noise, decreasing density at
${\sim}1\%$/month~\cite{Lang2004}.
This is not pathology---it is C2 operating correctly
with zero input.
\end{remark}


% ------ 3.3 VASCULAR: GLAGOV ------
\subsection{Vascular: Glagov's Remodeling (1987)}
\label{sec:glagov}

Glagov observed that arteries compensatorily enlarge
as atherosclerotic plaque accumulates, maintaining
lumen diameter until plaque burden exceeds
${\sim}40\%$~\cite{Glagov1987}.

The mechanotransduction pathway is:
wall shear stress ($\tau_w$) $\to$ endothelial
nitric oxide $\to$ smooth muscle relaxation $\to$
vessel diameter increase $\to$ reduced $\tau_w$.

In impedance language:
\begin{equation}
  Z_{\text{vasc}} = \frac{\Delta P}{Q}
  = \frac{8\mu L}{\pi r^4}
  \quad \text{(Poiseuille)}\,,
\end{equation}
where $r$ is lumen radius.

The remodeling equation is:
\begin{equation}
  \frac{dr}{dt} = k(\tau_w - \tau_{\text{ref}})\,,
  \label{eq:glagov-shear}
\end{equation}
which changes $Z_{\text{vasc}} \propto r^{-4}$.

Since $\tau_w \propto Q/r^3$ is the flow-derived signal
arriving at the wall ($x_{\text{pre}}$), and the
endothelial response is the transmitted signal
($x_{\text{post}}$), and the mismatch
$\tau_w - \tau_{\text{ref}}$ maps to $\Gamma$:
\begin{equation}
  \Delta Z_{\text{vasc}}
  = -\eta\,\Gamma_{\text{vasc}}\,
    \tau_w\,(\text{endothelial response})
  = -\eta\,\Gamma\,x_{\text{pre}}\,x_{\text{post}}\,.
  \quad \checkmark
  \label{eq:glagov-C2}
\end{equation}

\begin{remark}
The $40\%$ threshold in Glagov's data corresponds to
the critical $\Gamma$ value beyond which the compensatory
remodeling rate (limited by smooth muscle proliferation)
cannot keep pace with plaque growth---i.e., the learning
rate $\eta$ is too small relative to the impedance
perturbation rate.
\end{remark}


% ------ 3.4 IMMUNE: MEMORY ------
\subsection{Immune: Affinity Maturation and Memory}
\label{sec:immune}

Upon antigen exposure, B cells undergo somatic hypermutation
in germinal centres, producing variants with higher affinity
($K_a$) for the antigen~\cite{Victora2012}.

In impedance language, the receptor--antigen interaction
impedance is:
\begin{equation}
  Z_{\text{receptor}} = \frac{1}{K_a \cdot [\text{receptor}]}\,,
\end{equation}
where low $Z$ means efficient binding (good impedance match).

The first infection: $Z_{\text{immune}}$ is high
(na\"ive; poor match to pathogen) $\Rightarrow$ large
$\Gamma$ $\Rightarrow$ reflected energy (inflammation,
fever, tissue damage = wasted energy).

Affinity maturation is:
\begin{equation}
  \Delta K_a \propto K_a \cdot [\text{Ag}] \cdot [\text{T}_{\text{FH}}]\,,
\end{equation}
where $[\text{Ag}]$ is antigen concentration
($x_{\text{pre}}$) and $[\text{T}_{\text{FH}}]$ is
T follicular helper signal ($x_{\text{post}}$).
Since $\Delta Z \propto -\Delta K_a$:
\begin{equation}
  \Delta Z_{\text{receptor}}
  = -\eta\,\Gamma_{\text{immune}}\,
    [\text{Ag}]\,[\text{T}_{\text{FH}}]
  = -\eta\,\Gamma\,x_{\text{pre}}\,x_{\text{post}}\,.
  \quad \checkmark
  \label{eq:immune-C2}
\end{equation}

\begin{remark}
A \emph{vaccine} is the deliberate injection of
$x_{\text{pre}}$ (attenuated antigen) to drive the Hebbian
update $\Delta Z$ in the absence of disease.
Booster doses are repeated stimulation to further reduce
$\Gamma$, exactly as spaced repetition strengthens synaptic
potentiation~\cite{Cepeda2006}.
\end{remark}


% ------ 3.5 EPITHELIAL: KERATINIZATION ------
\subsection{Epithelial: Skin Keratinization}
\label{sec:skin}

Chronic mechanical friction on skin triggers accelerated
keratinocyte proliferation and cornification, forming
a callus---a region of increased mechanical
impedance~\cite{Thomas2010}.

The stimulus is external friction ($x_{\text{pre}}$;
mechanical force arriving at the skin interface).
The response is tissue deformation ($x_{\text{post}}$;
the transmitted mechanical signal).
The mismatch between the hard tool and the soft skin is:
\begin{equation}
  \Gamma_{\text{skin}} =
  \frac{Z_{\text{tool}} - Z_{\text{skin}}}
       {Z_{\text{tool}} + Z_{\text{skin}}}\,.
\end{equation}
Keratinization increases $Z_{\text{skin}}$ toward
$Z_{\text{tool}}$:
\begin{equation}
  \Delta Z_{\text{skin}}
  = -\eta\,\Gamma_{\text{skin}}\,
    F_{\text{friction}}\,\delta_{\text{tissue}}
  = -\eta\,\Gamma\,x_{\text{pre}}\,x_{\text{post}}\,.
  \quad \checkmark
  \label{eq:skin-C2}
\end{equation}
Note $\Gamma_{\text{skin}} < 0$ (soft skin vs.\ hard tool),
so $\Delta Z > 0$: impedance increases.
When the callus hardness $\approx$ tool hardness,
$\Gamma \to 0$ and remodeling stops.

Cessation of friction ($x_{\text{pre}} = 0$) halts the
Hebbian drive, and keratinocyte turnover reverts
$Z_{\text{skin}}$ to baseline---the callus
disappears~\cite{Thomas2010}.

\begin{remark}
The same mechanism applies to vocal cord nodules in
singers (repeated vibratory friction) and to the
thickened plantar fascia in runners.
In every case, the tissue impedance adapts toward
matching the source impedance of the applied load.
\end{remark}


% ============================================================
%  SECTION 3.6–3.10 — REMAINING FIVE NETWORKS (FULL PROOFS)
% ============================================================
\subsection{Muscular: Davis's Law (1867) and Hill's Equation}
\label{sec:muscle}

Davis's Law~\cite{Davis1867} states that soft tissue
remodels along lines of imposed mechanical demand.
Muscle hypertrophy and atrophy are the most visible
examples of tissue-level impedance adaptation in
daily life.

\paragraph{Impedance definition.}
Muscle mechanical impedance is the ratio of contractile
force to shortening velocity---precisely Hill's
equation~\cite{Hill1938}:
\begin{equation}
  Z_{\text{muscle}} = \frac{F}{v}
  = \frac{(F_0 + a)(v + b)}{v} - a\,,
  \label{eq:hill}
\end{equation}
where $F_0$ is the isometric force related to
cross-sectional area (CSA), and $a$, $b$ are
Hill's constants determined by muscle fiber type
and temperature.  For the impedance-matching argument,
the key insight is that $Z_{\text{muscle}}$ increases
with CSA (more sarcomeres in parallel $\Rightarrow$
higher force output at given velocity).

\paragraph{The muscular $\Gamma$.}
\begin{equation}
  \Gamma_{\text{muscle}} =
  \frac{Z_{\text{load}} - Z_{\text{muscle}}}
       {Z_{\text{load}} + Z_{\text{muscle}}}\,,
  \label{eq:gamma-muscle}
\end{equation}
where $Z_{\text{load}}$ is the mechanical impedance
of the external load (object weight, gravitational
demand, or antagonist tension).

\paragraph{C2 mapping.}
Assign $x_{\text{pre}} = $ neural drive from
$\alpha$-motor neurons (recruitment level) and
$x_{\text{post}} = $ muscle tension (the mechanical
response of the fiber to activation):
\begin{equation}
  \Delta Z_{\text{muscle}} =
  -\eta\,\Gamma_{\text{muscle}}\,
  (\text{neural drive})\,(\text{muscle tension})\,.
  \label{eq:c2-muscle}
\end{equation}

\paragraph{Mechanism.}
When a muscle is repeatedly loaded beyond its matched
range ($\Gamma > 0$, load exceeds muscle capacity),
the Hebbian product (neural activation $\times$
tension) is large, driving $\Delta Z > 0$
(CSA increases via satellite cell recruitment,
myofibrillar addition).
This is \textbf{hypertrophy}: the muscle adapts
its impedance upward to match the chronic load.

When loading is removed ($x_{\text{pre}} \to 0$,
bed rest or spinal cord injury), the Hebbian
product vanishes, C2 produces no update, and
baseline protein degradation causes CSA to decline.
This is \textbf{atrophy}---not a separate mechanism
but C2 with $x_{\text{pre}} = 0$.

\paragraph{Pathology: sarcopenia.}
Age-related sarcopenia involves declining motor
neuron recruitment ($x_{\text{pre}} \downarrow$)
and reduced satellite cell pool ($\eta \downarrow$).
Both reduce C2 drive, causing progressive
impedance mismatch with even normal gravitational
loads.  The result is $\Gamma_{\text{muscle}} > 0$
that the aging tissue cannot correct.

\paragraph{Training is C2 optimization.}
Resistance training protocols are empirically
optimized to maximize the $x_{\text{pre}} \cdot
x_{\text{post}}$ product: heavy loads (high tension)
with full motor recruitment (high neural drive)
produce the largest $\Delta Z$ per session.
This is why low-load/high-repetition and
high-load/low-repetition training produce
equivalent hypertrophy at matched
volume~\cite{Schoenfeld2017}---both maximize
the integrated Hebbian product over different
time profiles.


\subsection{Pulmonary: Altitude Acclimatization and Respiratory Remodeling}
\label{sec:pulmonary}

\paragraph{Impedance definition.}
Airway impedance is the ratio of airway pressure
to volumetric gas flow:
\begin{equation}
  Z_{\text{aw}} = \frac{P_{\text{aw}}}{\dot{V}}\,,
  \label{eq:z-aw}
\end{equation}
which has a resistive component (airway caliber)
and a compliant component (alveolar elastance).
Total pulmonary impedance includes the gas-exchange
interface: alveolar surface area $A$ determines
diffusion conductance $D_L = A \cdot D_m / T$
(Fick's law), and $Z_{\text{gas}} \propto 1/D_L$.

\paragraph{The pulmonary $\Gamma$.}
\begin{equation}
  \Gamma_{\text{pulm}} =
  \frac{Z_{\text{demand}} - Z_{\text{lung}}}
       {Z_{\text{demand}} + Z_{\text{lung}}}\,,
  \label{eq:gamma-pulm}
\end{equation}
where $Z_{\text{demand}}$ is set by metabolic
oxygen requirement and $Z_{\text{lung}}$ is the
lung's current capacity.

\paragraph{C2 mapping.}
\begin{equation}
  \Delta Z_{\text{lung}} =
  -\eta\,\Gamma_{\text{pulm}}\,
  (\text{ventilatory drive})\,
  (\text{alveolar expansion})\,.
  \label{eq:c2-pulm}
\end{equation}

\paragraph{Altitude acclimatization.}
At high altitude ($>3000$~m), hypoxia increases
ventilatory drive ($x_{\text{pre}} \uparrow$).
Alveolar expansion ($x_{\text{post}}$) is maximal
during hyperpnea.  Over days to weeks, C2 drives:
(1)~pulmonary vascular remodeling (reducing pulmonary
vascular resistance), (2)~erythropoietin-driven
polycythemia (increasing O$_2$ carrying capacity),
and (3)~increased alveolar surface via
capillary proliferation~\cite{Hackett2001,West2012}.
All three reduce $Z_{\text{lung}}$ toward
$Z_{\text{demand}}$.  Acclimatization is complete
when $\Gamma_{\text{pulm}} \to 0$ (gas exchange
matches metabolic demand at altitude).

\paragraph{COPD: irreversible $\Gamma$.}
Chronic obstructive pulmonary disease
destroys alveolar architecture (emphysema).
Unlike altitude acclimatization, the alveolar damage
is \emph{irreversible}---the tissue cannot execute
C2 because the substrate for remodeling has been
ablated.  The result is permanently elevated
$\Gamma_{\text{pulm}}$:
the lung cannot match even resting metabolic
demand, and the compensatory increase in
ventilatory drive ($x_{\text{pre}}$) provides
only partial relief via airway remodeling
in surviving tissue.

\paragraph{Infant lung development.}
Neonatal lungs have small alveolar surface area
($Z_{\text{lung}} \gg Z_{\text{demand}}$,
$\Gamma_{\text{pulm}} \gg 0$).
The first cry generates massive ventilatory drive;
over weeks to months, alveolar multiplication
progressively reduces $Z$.  By age 2--3,
$\Gamma_{\text{pulm}} \approx 0$ for resting
metabolism.  This is C2 operating at developmental
timescale ($\eta_{\text{dev}} \gg
\eta_{\text{acute}}$).


\subsection{Gastrointestinal: Intestinal Adaptation and the Gut-Brain Axis}
\label{sec:GI}

\paragraph{Impedance definition.}
Absorptive impedance is the ratio of luminal
nutrient concentration to transmucosal absorption
flux:
\begin{equation}
  Z_{\text{GI}} = \frac{[\text{nutrient}]_{\text{lumen}}}
                       {J_{\text{absorption}}}\,,
  \label{eq:z-gi}
\end{equation}
which depends on villous surface area, transporter
density, and mucosal blood flow.  Higher surface area
$\Rightarrow$ lower $Z_{\text{GI}}$.

\paragraph{The intestinal $\Gamma$.}
\begin{equation}
  \Gamma_{\text{GI}} =
  \frac{Z_{\text{dietary}} - Z_{\text{GI}}}
       {Z_{\text{dietary}} + Z_{\text{GI}}}\,,
  \label{eq:gamma-gi}
\end{equation}
where $Z_{\text{dietary}}$ represents the impedance
of the dietary load (food composition, volume,
frequency).  Simple carbohydrates have low
$Z_{\text{dietary}}$ (easy to absorb);
complex proteins have high $Z_{\text{dietary}}$
(require extensive enzymatic processing).

\paragraph{C2 mapping.}
\begin{equation}
  \Delta Z_{\text{GI}} =
  -\eta\,\Gamma_{\text{GI}}\,
  (\text{luminal content})\,
  (\text{epithelial uptake})\,.
  \label{eq:c2-gi}
\end{equation}

\paragraph{Post-resection adaptation.}
After partial bowel resection, the absorptive
surface area is abruptly reduced ($Z_{\text{GI}}
\uparrow$, $\Gamma \gg 0$).
Over weeks to months, the remaining intestine
undergoes compensatory hypertrophy:
villous height increases 50--100\%, crypt depth
doubles, and transporter expression
upregulates~\cite{Tappenden2014}.
All three reduce $Z_{\text{GI}}$ until
$\Gamma \to 0$ and nutritional homeostasis is
restored.  This is textbook C2: the tissue
adapts its impedance to match the increased
nutrient load per unit length.

\paragraph{The enteric nervous system.}
The gut contains $\sim 500$ million neurons
(the enteric nervous system, ENS)---a
complete $\Gamma$-subnetwork that operates
quasi-autonomously.  The ENS implements local C2:
sensory neurons detect luminal $Z_{\text{dietary}}$,
motor neurons regulate peristalsis and secretion,
and synaptic plasticity adjusts the
enteric circuit to dietary patterns.
The vagus nerve couples this subnetwork to the
brainstem controller (Section~\ref{sec:regulatory-brain}),
enabling top-down modulation of gut impedance
via the gut--brain axis.

\paragraph{Pathology: celiac disease.}
In celiac disease, gluten triggers autoimmune
villous atrophy---$Z_{\text{GI}}$ rises
because the absorptive surface is destroyed.
Unlike post-resection adaptation, the destructive
stimulus is \emph{ongoing} (continued gluten exposure),
and the immune attack prevents C2 from executing
compensatory growth.  The treatment (gluten-free diet)
removes the destructive signal, allowing C2 to
rebuild villi over months.


\subsection{Endocrine: Receptor Regulation and Metabolic Impedance}
\label{sec:endocrine}

\paragraph{Impedance definition.}
Endocrine impedance is defined through receptor
binding kinetics:
\begin{equation}
  Z_{\text{endo}} = \frac{1}{K_d \cdot N_R}\,,
  \label{eq:z-endo}
\end{equation}
where $K_d$ is the receptor dissociation constant
and $N_R$ is receptor density.
Low $Z_{\text{endo}}$ means high sensitivity
(many high-affinity receptors);
high $Z_{\text{endo}}$ means low sensitivity.

\paragraph{The endocrine $\Gamma$.}
\begin{equation}
  \Gamma_{\text{endo}} =
  \frac{Z_{\text{hormone}} - Z_{\text{endo}}}
       {Z_{\text{hormone}} + Z_{\text{endo}}}\,,
  \label{eq:gamma-endo}
\end{equation}
where $Z_{\text{hormone}} \propto 1/[\text{H}]$
(inversely proportional to hormone concentration).
When hormone is in excess, $Z_{\text{hormone}}$
is low, creating $\Gamma < 0$.

\paragraph{C2 mapping.}
\begin{equation}
  \Delta Z_{\text{endo}} =
  -\eta\,\Gamma_{\text{endo}}\,
  (\text{hormone conc.})\,
  (\text{intracellular response})\,.
  \label{eq:c2-endo}
\end{equation}

\paragraph{Bidirectional regulation.}
\emph{Downregulation} (chronic hormone excess):
$\Gamma < 0$ drives $\Delta Z > 0$
(receptors internalize, $N_R$ decreases,
$K_d$ increases).
\emph{Upregulation} (chronic deficiency):
$\Gamma > 0$ drives $\Delta Z < 0$
(receptor synthesis increases, sensitivity rises).
Both are the same C2 equation applied with
opposite $\Gamma$ signs.

This symmetry is not present in any single
biochemical signaling model---it emerges
naturally from the impedance-matching framework
because $\Gamma$ has a sign that automatically
selects the correct adaptation direction.

\paragraph{Insulin resistance: $\Gamma$ saturation.}
Type~2 diabetes exemplifies C2 failure.
Chronic hyperinsulinemia keeps
$Z_{\text{hormone}}$ persistently low.
C2 correctly drives receptor downregulation
($Z_{\text{endo}} \uparrow$), but this
\emph{worsens} the mismatch in a runaway loop:
downregulated receptors cause the pancreas to
secrete more insulin, further lowering
$Z_{\text{hormone}}$~\cite{Kahn2006}.
Eventually $\Gamma_{\text{endo}}$ saturates at
a large negative value, and
$\Delta Z$ per tick cannot escape the positive
feedback trap.

This is formally analogous to Glagov's 40\%
threshold in vascular remodeling: beyond a
critical $|\Gamma|$, C2's linear learning rate
is insufficient to correct an exponentially
growing stimulus.

\paragraph{Thyroid as calibration organ.}
Thyroid hormone (T$_3$/T$_4$) sets the metabolic
rate of virtually every cell, making it a
``gain control'' for the entire impedance network.
Hyperthyroidism ($Z_{\text{hormone}}$ too low)
drives global receptor downregulation;
hypothyroidism drives upregulation.
The endocrine $\Gamma$-network thus acts as a
global scaling factor for all other networks'
$\eta$ values---a meta-controller
distinct from both the cardiac pump
and the brainstem setpoint controller.


\subsection{Lymphatic: Reactive Hypertrophy and Immune Drainage}
\label{sec:lymphatic}

\paragraph{Impedance definition.}
Lymphatic drainage impedance is the ratio of
interstitial pressure to immune cell transport
flow:
\begin{equation}
  Z_{\text{lymph}} =
  \frac{P_{\text{interstitial}}}{Q_{\text{lymph}}}\,,
  \label{eq:z-lymph}
\end{equation}
where $Q_{\text{lymph}}$ is lymph flow rate
(determined by lymph vessel contractility,
node filtering capacity, and tissue drainage).
A larger, more vascularized lymph node has
lower $Z_{\text{lymph}}$ (more throughput
capacity per unit pressure).

\paragraph{The lymphatic $\Gamma$.}
\begin{equation}
  \Gamma_{\text{lymph}} =
  \frac{Z_{\text{load}} - Z_{\text{lymph}}}
       {Z_{\text{load}} + Z_{\text{lymph}}}\,,
  \label{eq:gamma-lymph}
\end{equation}
where $Z_{\text{load}}$ is the impedance of
the immune challenge (antigen load, tissue
inflammation, infection burden).
Higher pathogen load $\Rightarrow$ higher
$Z_{\text{load}}$ $\Rightarrow$ $\Gamma > 0$.

\paragraph{C2 mapping.}
\begin{equation}
  \Delta Z_{\text{lymph}} =
  -\eta\,\Gamma_{\text{lymph}}\,
  (\text{tissue antigen load})\,
  (\text{lymph node response})\,.
  \label{eq:c2-lymph}
\end{equation}

\paragraph{Reactive lymphadenopathy.}
During regional infection, antigen-presenting
cells activate lymph node follicular expansion
(germinal center reaction).
The $x_{\text{pre}} \cdot x_{\text{post}}$ product
is large (high antigen load $\times$ vigorous
immune response), driving $\Delta Z < 0$:
the lymph node enlarges, adds high-endothelial
venules (HEV), and increases filtration capacity.
Clinically, this is \textbf{reactive
lymphadenopathy}---lymph nodes swell precisely
because they are executing C2 to match the
increased immune traffic demand.

When infection resolves ($x_{\text{pre}} \to 0$),
the Hebbian product vanishes, C2 drive ceases,
and the node involutes to baseline via
apoptosis of expanded clones.

\paragraph{Pathology: lymphedema.}
Lymphedema (post-surgical, post-radiation)
occurs when lymphatic vessels or nodes are
physically destroyed.
The remaining lymphatic network cannot execute C2
because the tissue substrate is ablated
($\eta \to 0$).  $\Gamma_{\text{lymph}}$ is
permanently elevated: even normal interstitial
fluid production exceeds drainage capacity,
causing chronic edema.

\paragraph{Sentinel node as impedance probe.}
In oncology, the sentinel lymph node is the
first node to receive drainage from a tumor.
Its impedance state reflects the tumor's
$Z_{\text{load}}$: early metastasis changes
the node's architecture (C2 driven by tumor
antigens), which is detectable by histology.
The sentinel node biopsy is, physically,
a measurement of $\Gamma_{\text{lymph}}$ at
the tumor-drainage interface.


\subsection{Renal: Tubuloglomerular Feedback}
\label{sec:renal}

The kidney is the body's impedance-matching engine for
blood composition: it filters, reabsorbs, and secretes
to maintain electrolyte, acid--base, and volume homeostasis.

\paragraph{Impedance definition.}
Renal hemodynamic impedance at the glomerulus is:
\begin{equation}
  Z_{\text{renal}} =
  \frac{P_{\text{glom}}}{Q_{\text{filtrate}}}\,,
  \label{eq:z-renal}
\end{equation}
where $P_{\text{glom}}$ is glomerular capillary pressure
and $Q_{\text{filtrate}}$ is the glomerular filtration
rate (GFR).  Afferent and efferent arteriolar tone
set $Z_{\text{renal}}$; tubular reabsorption capacity
sets $Z_{\text{tubule}}$.

\paragraph{The renal $\Gamma$.}
\begin{equation}
  \Gamma_{\text{renal}} =
  \frac{Z_{\text{load}} - Z_{\text{renal}}}
       {Z_{\text{load}} + Z_{\text{renal}}}\,,
  \label{eq:gamma-renal}
\end{equation}
where $Z_{\text{load}}$ is the filtration demand
imposed by plasma composition (solute load, volume).

\paragraph{C2 mapping.}
Tubuloglomerular feedback (TGF) senses NaCl
concentration at the macula densa
($x_{\text{pre}} = $ macula densa signal) and
adjusts afferent arteriolar tone
($x_{\text{post}} = $ vascular smooth muscle
response):
\begin{equation}
  \Delta Z_{\text{renal}} =
  -\eta\,\Gamma_{\text{renal}}\,
  (\text{macula densa signal})\,
  (\text{arteriolar response})\,.
  \label{eq:c2-renal}
\end{equation}

This is a negative feedback loop that stabilizes
single-nephron GFR within seconds---one of the fastest
C2 implementations in the body
($\eta_{\text{TGF}} \gg \eta_{\text{bone}}$).

\paragraph{Chronic kidney disease: nephron dropout.}
When nephrons are lost (diabetes, hypertension),
surviving nephrons undergo compensatory hyperfiltration:
$Z_{\text{renal}}$ per nephron decreases to maintain
total GFR.  This is C2 operating correctly.
However, the increased single-nephron $\Gamma$
(higher demand per nephron) eventually causes
glomerulosclerosis---the nephron's C2 capacity is
exceeded, creating a positive feedback loop
of nephron loss~\cite{Brenner1982}.
CKD progression is the renal analog of Glagov's
40\% threshold in vascular remodeling.

\paragraph{RAAS as second-level controller.}
The renin--angiotensin--aldosterone system operates
as a renal-specific controller
(analogous to the cardiac pump for the vascular tree):
it adjusts systemic $Z_{\text{renal}}$ via
angiotensin~II (vasoconstriction) and aldosterone
(Na$^+$ reabsorption).  ACE inhibitors work by
reducing $Z_{\text{renal}}$, relieving the
chronic $\Gamma$ that drives progressive damage.


\subsection{Hepatic: Liver Regeneration}
\label{sec:hepatic}

The liver's regenerative capacity---recovering full mass
after 70\% resection within weeks---is arguably the most
dramatic C2 response in any mammalian tissue.

\paragraph{Impedance definition.}
Hepatic metabolic impedance is the ratio of portal
blood composition (metabolite concentration) to
hepatic clearance flux:
\begin{equation}
  Z_{\text{hepatic}} =
  \frac{[\text{metabolite}]_{\text{portal}}}
       {J_{\text{clearance}}}\,.
  \label{eq:z-hepatic}
\end{equation}
After partial hepatectomy, $J_{\text{clearance}}$
drops abruptly (fewer hepatocytes), so
$Z_{\text{hepatic}}$ rises sharply.

\paragraph{The hepatic $\Gamma$.}
\begin{equation}
  \Gamma_{\text{hepatic}} =
  \frac{Z_{\text{metabolic\,load}} - Z_{\text{hepatic}}}
       {Z_{\text{metabolic\,load}} + Z_{\text{hepatic}}}\,.
  \label{eq:gamma-hepatic}
\end{equation}

\paragraph{C2 mapping.}
Portal blood delivers growth factors and metabolites
that act as the pre-synaptic signal
($x_{\text{pre}} = $ portal flow composition:
HGF, bile acids, cytokines);
the hepatocyte proliferative response is
$x_{\text{post}} = $ mitotic activity:
\begin{equation}
  \Delta Z_{\text{hepatic}} =
  -\eta\,\Gamma_{\text{hepatic}}\,
  (\text{portal signal})\,
  (\text{proliferative response})\,.
  \label{eq:c2-hepatic}
\end{equation}

After 70\% hepatectomy, $\Gamma \gg 0$ (metabolic
demand vastly exceeds clearance capacity).
The remaining hepatocytes enter rapid
proliferation~\cite{Michalopoulos2007},
driven by the massive $x_{\text{pre}} \cdot
x_{\text{post}}$ Hebbian product.
Regeneration halts precisely when liver mass
restores $Z_{\text{hepatic}}$ to match the
metabolic load ($\Gamma \to 0$)---not at any
fixed anatomical size, but at functional
impedance balance.

\paragraph{Termination signal is $\Gamma = 0$.}
This explains the long-standing puzzle of how
the liver ``knows'' when to stop regenerating:
it does not count cells or measure mass.
It measures $\Gamma_{\text{hepatic}}$ (the
mismatch between portal metabolite load and
clearance capacity).  When $\Gamma \to 0$,
the Hebbian product vanishes and proliferation
stops.

\paragraph{Cirrhosis: $\eta \to 0$.}
In cirrhosis, fibrosis replaces functional
hepatocytes with scar tissue.
The metabolic $\Gamma$ remains high, but
the regenerative learning rate $\eta$
collapses because the tissue substrate for
proliferation is destroyed.
$Z_{\text{hepatic}}$ cannot decrease,
$\Gamma$ remains permanently elevated,
and hepatic decompensation follows.


\subsection{Connective Tissue: Tendon and Ligament Remodeling}
\label{sec:connective}

Tendons and ligaments occupy the mechanical interface
between bone (Wolff) and muscle (Davis), transmitting
forces across joints.

\paragraph{Impedance definition.}
Tendon mechanical impedance is viscoelastic stiffness:
\begin{equation}
  Z_{\text{tendon}} = E \cdot \frac{A}{L}\,,
  \label{eq:z-tendon}
\end{equation}
where $E$ is Young's modulus (collagen fiber
organization and cross-linking), $A$ is cross-sectional
area, and $L$ is rest length.

\paragraph{C2 mapping.}
\begin{equation}
  \Delta Z_{\text{tendon}} =
  -\eta\,\Gamma_{\text{tendon}}\,
  (\text{mechanical strain})\,
  (\text{fibroblast response})\,.
  \label{eq:c2-tendon}
\end{equation}
Tenocytes (fibroblasts) sense mechanical strain via
integrin-mediated mechanotransduction ($x_{\text{pre}}$)
and respond with collagen synthesis and
alignment ($x_{\text{post}}$).
Repeated loading aligns collagen fibers along
principal stress directions and increases
cross-sectional area---both reduce $\Gamma$ by
increasing $Z_{\text{tendon}}$ toward the
load impedance.

\paragraph{Timescale hierarchy.}
Connective tissue has an intermediate $\eta$:
slower than muscle ($\sim$weeks vs.\ days) but
faster than bone ($\sim$weeks vs.\ months).
This creates a temporal impedance-matching chain:
\emph{muscle adapts first}, then \emph{tendon},
then \emph{bone}.  Training injuries often occur
at the tendon because muscle C2 outpaces
tendon C2, creating a transient $\Gamma$ spike
at the myotendinous junction.

\paragraph{Pathology: tendinopathy.}
Chronic overuse without adequate recovery maintains
$\Gamma > 0$ while the repair cycle
(collagen turnover) cannot keep pace.
The tendon enters a degenerative state:
disorganized collagen, neovascularization,
increased $Z$ variability.
This is not inflammation but failed C2---the
tissue is attempting impedance matching but
the stimulus exceeds $\eta$.


\subsection{Adipose: Energy Storage and Leptin Signaling}
\label{sec:adipose}

Adipose tissue is not passive storage; it is the body's
energy impedance buffer and an endocrine organ.

\paragraph{Impedance definition.}
Adipose metabolic impedance is the ratio of
circulating energy substrate to storage/release flux:
\begin{equation}
  Z_{\text{adipose}} =
  \frac{[\text{energy substrate}]}
       {J_{\text{storage/release}}}\,.
  \label{eq:z-adipose}
\end{equation}
In energy excess, substrate concentration rises and
storage flux increases ($Z$ adjusts via adipocyte
hypertrophy and hyperplasia).

\paragraph{C2 mapping.}
\begin{equation}
  \Delta Z_{\text{adipose}} =
  -\eta\,\Gamma_{\text{adipose}}\,
  (\text{insulin signal})\,
  (\text{adipocyte response})\,.
  \label{eq:c2-adipose}
\end{equation}
When caloric intake exceeds expenditure,
$\Gamma_{\text{adipose}} > 0$: C2 drives
adipocyte expansion (hypertrophy then hyperplasia)
to increase storage capacity, reducing $\Gamma$.

\paragraph{Leptin as impedance signal.}
Leptin secretion is proportional to adipose mass
(i.e., to $Z_{\text{adipose}}$).
It acts as an impedance report to the hypothalamic
controller (Section~\ref{sec:regulatory-brain}):
high leptin signals $\Gamma_{\text{adipose}} \approx 0$
(energy stores adequate);
low leptin signals $\Gamma > 0$
(energy deficit)~\cite{Friedman2019}.

\paragraph{Leptin resistance $\equiv$ insulin resistance.}
Obesity produces chronic hyperleptinemia.
The hypothalamic controller downregulates leptin
receptors ($Z_{\text{setpoint}}$ drifts upward),
analogous to insulin resistance in the endocrine
network (Section~\ref{sec:endocrine}).
The controller can no longer ``see'' the energy
surplus: $\Gamma_{\text{reg}}$ is miscalibrated,
and the organism behaves as though starving despite
excess stores.  This is a controller-level C2 failure,
not a tissue-level failure---the adipose C2 is
functioning correctly, but the regulatory brain
has lost impedance calibration.


\subsection{Hematopoietic: Erythropoiesis and Blood Cell Production}
\label{sec:hematopoietic}

The hematopoietic system continuously produces $\sim 200$
billion red blood cells per day, adjusting production rate
to match oxygen delivery demand.

\paragraph{Impedance definition.}
Oxygen transport impedance:
\begin{equation}
  Z_{\text{hema}} =
  \frac{P_{O_2,\text{demand}}}
       {Q_{O_2,\text{delivery}}}\,,
  \label{eq:z-hema}
\end{equation}
where delivery depends on hemoglobin concentration,
cardiac output, and SpO$_2$.

\paragraph{C2 mapping.}
\begin{equation}
  \Delta Z_{\text{hema}} =
  -\eta\,\Gamma_{\text{hema}}\,
  (\text{EPO signal})\,
  (\text{erythroid response})\,.
  \label{eq:c2-hema}
\end{equation}
Renal hypoxia sensors detect oxygen deficit
($\Gamma_{\text{hema}} > 0$) and secrete
erythropoietin (EPO), which is $x_{\text{pre}}$.
Bone marrow erythroid progenitors respond
by accelerating red cell production
($x_{\text{post}}$), increasing hemoglobin
concentration and thus $Q_{O_2}$,
driving $\Gamma \to 0$.

\paragraph{Altitude erythropoiesis.}
At high altitude, $\Gamma_{\text{hema}} > 0$
(demand exceeds delivery at reduced $P_{O_2}$).
EPO rises 2--3 fold; reticulocyte count
peaks at $\sim 7$~days.  Over weeks,
hemoglobin rises until $\Gamma \to 0$ at the
new altitude~\cite{Hackett2001}.
This connects to the pulmonary $\Gamma$-network
(Section~\ref{sec:pulmonary}): both adapt
\emph{simultaneously} to the same hypoxic stimulus
through different C2 pathways.

\paragraph{Anemia of chronic disease.}
Inflammatory cytokines suppress EPO production
and block iron release (hepcidin pathway).
The hematopoietic C2 is correct---$\Gamma > 0$
should drive increased production---but
$x_{\text{pre}}$ (EPO) is suppressed and iron
availability reduces $\eta$.
The system is ``trying'' to execute C2 but is
blocked by the immune network's cross-talk.

\paragraph{Polycythemia vera: C2 without $\Gamma$.}
In polycythemia vera, a JAK2 mutation drives
constitutive erythroid proliferation
regardless of $\Gamma$.
This is the hematopoietic equivalent of
cancer---C2 machinery running with
$x_{\text{pre}}$ permanently ``on'' despite
$\Gamma \approx 0$, producing dangerous
excess.


\subsection{Cartilage: The C2 Null Case ($\eta \approx 0$)}
\label{sec:cartilage}

Articular cartilage is scientifically important
not because it demonstrates C2, but because it
demonstrates what happens when C2 \emph{cannot execute}.

\paragraph{Impedance definition.}
Cartilage mechanical impedance (compressive stiffness):
\begin{equation}
  Z_{\text{cart}} = H_A \cdot \frac{A}{h}\,,
  \label{eq:z-cart}
\end{equation}
where $H_A$ is the aggregate modulus (a function of
proteoglycan content and collagen architecture) and
$h$ is cartilage thickness.

\paragraph{Why $\eta \approx 0$.}
Mature articular cartilage is avascular---it receives
nutrients solely by diffusion from synovial
fluid~\cite{Buckwalter2004}.
Chondrocytes are sparse ($< 5\%$ of tissue volume)
and post-mitotic in adults.
Consequently, the tissue's ability to remodel its
collagen network in response to mechanical mismatch
is negligible: $\eta_{\text{cart}} \approx 0$.

\paragraph{C2 predicts osteoarthritis.}
When joint loading changes (obesity, injury,
malalignment), $\Gamma_{\text{cart}} > 0$.
In any other tissue, C2 would drive impedance
adaptation.  In cartilage, $\eta \approx 0$
means $\Delta Z \approx 0$ regardless of $\Gamma$.
The tissue cannot adapt, and progressive
wear ensues.  Osteoarthritis is not a ``disease''
in the C2 framework---it is the \emph{predicted
consequence} of an interface with $\eta \approx 0$.

\paragraph{Theoretical significance.}
Cartilage is the \emph{falsifiable control case}
for the C2 universality claim:
if C2 is the universal remodeling law, then
tissues with $\eta \approx 0$ should show
irreversible degradation under chronic
$\Gamma > 0$.  Osteoarthritis confirms this
prediction exactly.
Conversely, tissue engineering efforts that
succeed in increasing chondrocyte proliferation
(raising $\eta$) should restore C2 remodeling
and halt OA progression---a testable prediction.


\subsection{Reproductive: Follicular Selection and Gametogenesis}
\label{sec:reproductive}

The reproductive system implements a unique form of C2:
\emph{competitive selection} among multiple candidate cells,
where only the best impedance-matched cell survives.

\paragraph{Ovarian follicular selection.}
Each menstrual cycle, multiple follicles compete for
FSH stimulation.  Define follicular impedance
as responsiveness to gonadotropin:
\begin{equation}
  Z_{\text{follicle}} = \frac{1}{N_{\text{FSH-R}} \cdot K_a}\,,
  \label{eq:z-follicle}
\end{equation}
where $N_{\text{FSH-R}}$ is FSH receptor density and
$K_a$ is binding affinity.  The follicle with the
lowest $Z$ (highest sensitivity) has the smallest
$\Gamma$ at the FSH interface and receives the
most signal energy:
\begin{equation}
  T_i = 1 - \Gamma_i^2 \propto
  \frac{1}{(1 + Z_{\text{follicle},i}/Z_{\text{FSH}})^2}\,.
  \label{eq:follicle-selection}
\end{equation}
As FSH levels decline (negative feedback from
the dominant follicle's estradiol), only the
best-matched follicle maintains sufficient
$T$ to survive.  Others undergo atresia
($T \to 0$, complete signal reflection).

\paragraph{C2 in follicular maturation.}
Each follicle executes C2 to improve its own
impedance matching:
\begin{equation}
  \Delta Z_{\text{follicle}} =
  -\eta\,\Gamma\,(\text{FSH})\,
  (\text{granulosa response})\,.
  \label{eq:c2-follicle}
\end{equation}
The follicle with the highest $\eta$ (most
responsive granulosa cells) reduces its $\Gamma$
fastest, becoming the dominant follicle.
This is a Darwinian competition implemented
through C2 at the receptor interface.

\paragraph{Spermatogenesis: temperature-dependent $\eta$.}
Spermatogenesis requires $\sim 34$\textdegree C
(below core body temperature).
In impedance terms, temperature sets the C2
learning rate $\eta$: at $37$\textdegree C
(cryptorchidism), $\eta \to 0$ and
spermatogenic C2 fails, causing infertility.
Testicular descent is the evolutionary solution
that places the tissue in the correct $\eta$ regime.

\paragraph{Pathology: PCOS.}
Polycystic ovary syndrome involves elevated
androgens that disrupt follicular C2: multiple
follicles begin maturation but none achieves
sufficient $\Gamma$ reduction for dominance.
The result is arrested follicles (``cysts'')---each
executing incomplete C2 without a winner.


% ============================================================
%  SECTION 3.18–3.19 — TWO CENTRAL CONTROLLERS
% ============================================================

\bigskip
\noindent
The preceding seventeen networks describe \emph{distributed}
tissue adaptation---local impedance matching along network
branches.  The next two describe the \emph{central
controllers} that set $Z_{\text{source}}$ or
$Z_{\text{setpoint}}$ for entire network trees.

\subsection{Cardiac: Laplace--Starling Remodeling}
\label{sec:cardiac}

The heart is not a branch of the vascular network; it is
the \emph{voltage source} (pump) whose output impedance
$Z_{\text{heart}}$ determines the pressure waveform
delivered to the entire arterial tree.
The aortic valve is the critical impedance interface:
\begin{equation}
  \Gamma_{\text{aortic}} =
  \frac{Z_{\text{aorta}} - Z_{\text{heart}}}
       {Z_{\text{aorta}} + Z_{\text{heart}}}\,.
  \label{eq:gamma-aortic}
\end{equation}

\paragraph{Impedance definitions.}
Ventricular output impedance is
$Z_{\text{heart}} = P_{\text{vent}} / Q_{\text{ej}}$
(ventricular pressure divided by ejection flow rate).
Aortic input impedance $Z_{\text{aorta}}$ is the
characteristic impedance of the aortic tree,
determined by vessel compliance and peripheral
resistance.

\paragraph{Laplace's law.}
Ventricular wall stress obeys:
\begin{equation}
  \sigma = \frac{P \cdot r}{2 h}\,,
  \label{eq:laplace}
\end{equation}
where $P$ is intraventricular pressure, $r$ is
chamber radius, and $h$ is wall thickness.

\paragraph{C2 mapping.}
Assign $x_{\text{pre}} = \sigma$ (the mechanical
stimulus sensed by cardiomyocytes) and
$x_{\text{post}} = \varepsilon$ (the myocardial strain
response).  Then C2 reads:
\begin{equation}
  \Delta Z_{\text{heart}} =
  -\eta\,\Gamma_{\text{aortic}}\,\sigma\,\varepsilon\,.
  \label{eq:c2-cardiac}
\end{equation}

\paragraph{Two modes of hypertrophy.}
When $Z_{\text{aorta}} > Z_{\text{heart}}$
(chronic pressure overload, e.g.\ aortic stenosis
or hypertension): $\Gamma > 0 \Rightarrow \Delta Z > 0
\Rightarrow$ wall thickness $h$ increases
$\Rightarrow \sigma$ decreases (Eq.~\ref{eq:laplace})
$\Rightarrow Z_{\text{heart}}$ rises toward
$Z_{\text{aorta}}$.
This is \textbf{concentric hypertrophy}.

When preload increases (chronic volume overload,
e.g.\ aortic regurgitation): the chamber dilates
($r \uparrow$), reducing the transmission
mismatch from a different direction.
This is \textbf{eccentric hypertrophy}.

Both are C2 with different boundary conditions.

\paragraph{Heart failure as $\Gamma \to 1$.}
Decompensated heart failure occurs when
$\Gamma_{\text{aortic}}$ is persistently elevated
and the C2 learning rate $\eta$ (cardiomyocyte
proliferation capacity) is insufficient to keep
pace~\cite{Grossman1975}.
The pump's output impedance remains mismatched,
causing chronic energy reflection into the
pulmonary circulation---the physical basis of
congestion.

\paragraph{Athlete's heart as healthy C2 cycling.}
Physiological cardiac hypertrophy in
athletes~\cite{Maron2006} demonstrates C2 under
favorable conditions: $\Gamma$ is periodically
elevated during exertion, C2 drives compensatory
remodeling, and rest periods allow recovery.
The result is balanced hypertrophy that improves
pump efficiency rather than causing failure.


\subsection{Regulatory Brain: Allostatic Adaptation}
\label{sec:regulatory-brain}

The brainstem (and hypothalamus) is not merely a neuron
collection that obeys Hebb's Law at its synapses.
It is the \emph{regulatory controller}
that sets the \emph{reference impedance}
$Z_{\text{setpoint}}$ for every homeostatic loop in
the body---blood pressure, heart rate, temperature,
cortisol, respiratory drive.

The controller's own setpoint drifts over time, and
this drift \emph{also} follows C2.

\paragraph{The regulatory reflection coefficient.}
Define the controller's error signal as:
\begin{equation}
  \Gamma_{\text{reg}} =
  \frac{Z_{\text{actual}} - Z_{\text{setpoint}}}
       {Z_{\text{actual}} + Z_{\text{setpoint}}}\,,
  \label{eq:gamma-reg}
\end{equation}
where $Z_{\text{actual}}$ is the measured peripheral
impedance (e.g.\ arterial pressure as seen by the
baroreceptor) and $Z_{\text{setpoint}}$ is the
brainstem's current reference value.

\paragraph{C2 mapping.}
Assign $x_{\text{pre}} = $ sensory afferent
(baroreceptor, chemoreceptor, thermoreceptor)
and $x_{\text{post}} = $ motor efferent
(sympathetic / parasympathetic output).  Then:
\begin{equation}
  \Delta Z_{\text{setpoint}} =
  -\eta\,\Gamma_{\text{reg}}\,
  x_{\text{afferent}}\,x_{\text{efferent}}\,.
  \label{eq:c2-regulatory}
\end{equation}

\paragraph{Baroreceptor resetting.}
This is exactly the clinically observed phenomenon of
baroreceptor resetting~\cite{Chapleau1991}:
chronic hypertension ($Z_{\text{actual}} >
Z_{\text{setpoint}}$, $\Gamma > 0$) drives the
setpoint upward ($Z_{\text{setpoint}} \uparrow$)
until $\Gamma \to 0$ at the new, elevated ``normal.''
The timescale is days to weeks---the $\eta$ of
brainstem plasticity.

\paragraph{HPA axis and allostatic load.}
The hypothalamic--pituitary--adrenal (HPA) axis
follows the same pattern.  Chronic stress
maintains $\Gamma_{\text{HPA}} > 0$;
C2 drives $Z_{\text{cortisol\_setpoint}}$ upward,
resulting in chronically elevated baseline cortisol.
McEwen~\cite{McEwen1998} termed the cumulative cost
of this setpoint drift \emph{allostatic load}.
In impedance language:
\begin{equation}
  \text{Allostatic load} =
  \int_0^T |\Delta Z_{\text{setpoint}}(t)|\,dt\,,
  \label{eq:allostatic}
\end{equation}
the total impedance drift accumulated by the
controller.  When this integral exceeds the
controller's recovery capacity, the organism
develops chronic disease---\emph{not} because
any tissue failed, but because the \emph{controller}
has drifted too far from physiological reference.

\paragraph{Why controllers are special.}
A failure in Glagov's vascular remodeling affects one
arterial branch.  A failure in baroreceptor resetting
or HPA adaptation affects \emph{every} organ simultaneously,
because $Z_{\text{setpoint}}$ propagates through the
entire autonomic tree.  This explains why chronic stress
causes multi-organ disease: it is not that stress
``damages everything''---it is that the \emph{controller's}
$\Gamma$ error propagates to all downstream networks.


% ============================================================
%  SECTION 4 — NHANES NETWORK PROPAGATION PROOF
% ============================================================
\section{Empirical Evidence: Network Propagation}
\label{sec:nhanes-proof}

\subsection{Experimental design}

Paper~VI validated $\Gamma$ diagnostics within single
organ systems.  Here we test a stronger prediction:
\emph{adding a physical measurement to one $\Gamma$-network
should improve predictions in coupled networks, through
impedance propagation at inter-network interfaces.}

Using NHANES 1999--2018 (10 cycles, $n = 49{,}774$,
$7{,}770$ deaths with up to 20-year follow-up), we
compared four configurations:
\begin{itemize}
  \item[A] Laboratory values only (27 items $\to$ 12 organs)
  \item[B] Labs $+$ blood pressure
    (SBP, DBP, pulse pressure $\to$ vascular, cardiac,
     renal, neurological $Z$)
  \item[C] Labs $+$ smoking status
    ($\to$ pulmonary, vascular, immune $Z$)
  \item[D] Full physics (Labs $+$ BP $+$ smoking)
\end{itemize}

All impedance mappings use textbook reference intervals
and organ weights derived from known pathophysiology.
\textbf{Zero parameters are fitted to any outcome data.}

\subsection{Results}

\subsubsection{All-cause mortality}

\begin{table}[h]
\centering
\caption{All-cause mortality AUC by configuration.
$\sum\Gamma^2$ is the total impedance mismatch burden.}
\label{tab:allcause}
\begin{tabular}{lcc}
\toprule
Configuration & AUC (95\% CI) & $\Delta$ vs.\ A \\
\midrule
A: Labs only    & 0.629 [0.622--0.636] & --- \\
B: Labs + BP    & 0.705 [0.699--0.711] & $+0.075$ \\
C: Labs + Smoke & 0.627 [0.620--0.634] & $-0.002$ \\
D: Full physics & 0.704 [0.698--0.710] & $+0.075$ \\
\bottomrule
\end{tabular}
\end{table}

The addition of blood pressure---a vascular network
measurement---increases the all-cause AUC by $12\%$
with zero fitted parameters.

\subsubsection{Network propagation: cross-organ improvement}

The critical test is whether a vascular input improves
non-vascular organ predictions.

\begin{table}[h]
\centering
\caption{Organ-specific death AUC: Labs only (A) vs.\
Labs + BP (B).  ``Network'' organs receive no direct BP
input; improvement comes from inter-network $\Gamma$-coupling.}
\label{tab:propagation}
\begin{tabular}{lcccc}
\toprule
Organ & Type & AUC (A) & AUC (B) & $\Delta$ \\
\midrule
Cardiac  & Network & 0.508 & 0.674 & $+0.166$ \\
Neuro    & Network & 0.516 & 0.577 & $+0.061$ \\
Renal    & Network & 0.677 & 0.700 & $+0.024$ \\
Endocrine & Control & 0.814 & 0.814 & $+0.000$ \\
\bottomrule
\end{tabular}
\end{table}

\textbf{Cardiac death AUC improves from 0.508 (chance) to
0.674}---a $\Delta$ of $+0.166$---from a vascular measurement,
not a cardiac measurement.
This occurs because SBP physically propagates through the
vascular--cardiac interface: the heart pumps against arterial
afterload, so vascular $Z$ directly determines cardiac $Z$.

Endocrine serves as a negative control: it shares no
significant interface with the vascular network for blood
pressure, and its AUC is unchanged ($\Delta = 0.000$).

\textbf{This cross-organ improvement is the empirical
signature of network propagation.}
The Framingham Risk Score also uses SBP,
but it produces a \emph{single} cardiovascular risk
estimate.  It cannot simultaneously improve renal and
neurological predictions because it has no organ network
topology.

\subsection{What the textbook gap means}

The difference between $\Gamma$ and textbook AUC is
not a model failure but an \emph{input measurement gap}:
\begin{equation}
  \text{AUC}_{\text{textbook}} - \text{AUC}_{\Gamma}
  = \text{unmapped tissue physics}\,.
  \label{eq:gap}
\end{equation}

For endocrine ($\Gamma$ AUC $= 0.814$), the physics
is nearly complete: glucose and HbA1c directly measure
endocrine impedance.
For cardiac ($\Gamma$ AUC $= 0.508 \to 0.674$),
adding SBP recovered 70\% of the gap to the SCORE2
benchmark (0.822).
The remaining gap requires troponin and BNP---cardiac
impedance measurements not available in NHANES.


% ============================================================
%  SECTION 5 — EVOLUTIONARY CONVERGENCE
% ============================================================
\section{Evolutionary Convergence on C2}
\label{sec:evolution}

\subsection{Why C2 is universal}

The gradient-descent derivation (Section~\ref{sec:c2-derivation})
shows that C2 is the \emph{unique} first-order
update rule that minimizes reflected energy at any
repeatedly stimulated interface.

Evolution ``discovers'' C2 not because organisms
``know'' calculus, but because organisms that happen
to encode C2-like chemistry at any interface
waste less energy, survive longer, and reproduce more.
The selection pressure is identical across all tissues:
reflected energy is wasted energy.

\subsection{The tissue is not the algorithm}

A key insight: the remodeling machinery differs vastly
across tissues.
\begin{itemize}
  \item Neural: NMDA receptors, CaMKII, AMPA trafficking.
  \item Bone: osteocyte mechanosensing, RANKL/OPG signaling.
  \item Vascular: endothelial shear sensors, eNOS, SMC
    proliferation.
  \item Immune: AID somatic hypermutation, T$_{\text{FH}}$
    selection.
  \item Skin: keratinocyte proliferation, calcium gradient
    signaling.
  \item Liver: HGF-driven hepatocyte mitosis.
  \item Kidney: macula densa NaCl sensing, afferent
    arteriolar tone.
  \item Adipose: insulin/leptin receptor signaling.
  \item Bone marrow: EPO-driven erythroid expansion.
  \item Reproductive: FSH receptor competition, granulosa
    proliferation.
\end{itemize}

These are ten completely different molecular machineries.
Yet they all implement the same equation because the
equation is dictated by physics, not by the substrate.
The substrate merely determines the material constants
($\eta$, the learning rate; $Z_0$, the reference impedance).

This is the biological analog of interface isomorphism
(Paper~VI, Proposition~1): the governing equation depends
on the \emph{boundary}, not on the bulk medium.

\subsection{No tissue is useless}

If $\text{Rank}(\mathbf{\Gamma}) = N_{\text{organs}}$
(every organ carries irreducible impedance information;
Paper~I), and every tissue remodels via C2, then
\emph{every tissue is an active participant in
minimizing $A[\Gamma]$}.

A tissue that appears ``useless'' in a clinical dataset
(e.g., cardiac AUC $= 0.508$) is not useless---it is
\emph{unobserved}.  The measurement matrix is rank-deficient,
not the biology.

Billions of years of evolution have not preserved any
tissue without function.  What we measure as
$\text{AUC} = 0.50$ is our blindness, not nature's waste.


% ============================================================
%  SECTION 6 — FALSIFIABLE PREDICTIONS
% ============================================================
\section{Falsifiable Predictions}
\label{sec:predictions}

The universal-C2 framework generates five testable predictions
that no single-organ model can produce:

\begin{enumerate}
  \item \textbf{Cross-organ propagation (confirmed).}
    Adding SBP (vascular) must improve cardiac, renal,
    and neurological AUC through network coupling.
    \emph{Result}: cardiac $+0.166$, neuro $+0.061$,
    renal $+0.024$. $\checkmark$

  \item \textbf{Control-organ invariance (confirmed).}
    Endocrine AUC must not change when SBP is added
    (no significant vascular--endocrine pressure interface).
    \emph{Result}: $\Delta = 0.000$. $\checkmark$

  \item \textbf{Microgravity bone loss follows C2 with
    $x_{\text{pre}} = 0$.}
    Bone density loss rate in microgravity should be
    predictable from $\eta$ and the pre-flight
    $\Gamma_{\text{mech}}$ distribution, without
    any fitted parameter.
    \emph{Status}: testable with ISS densitometry data.

  \item \textbf{Vaccine booster spacing is optimal when it
    maximizes $\sum x_{\text{pre}} \cdot x_{\text{post}}$.}
    The optimal booster interval should match the
    time at which memory B-cell $x_{\text{post}}$ signal
    is maximally available but $\Gamma_{\text{immune}}$
    has not yet fully decayed.
    \emph{Status}: testable against clinical immunization
    trial data.

  \item \textbf{Adding FEV1/SpO$_2$ (pulmonary $Z$) will
    improve immune death AUC.}
    The pulmonary--immune interface (lung-resident immune
    cells, mucosal barrier) should propagate pulmonary
    impedance information to immune predictions.
    \emph{Status}: testable with NHANES spirometry data
    (SPX files).
\end{enumerate}


% ============================================================
%  SECTION 7 — DISCUSSION
% ============================================================
\section{Discussion}
\label{sec:discussion}

\subsection{One equation, nineteen contexts}

The central message of this paper is not that Hebb's Law
``resembles'' Wolff's Law.  It is that both are
\emph{the same law}---the gradient-descent update
for minimizing reflected energy---expressed through
different molecular machineries.

Moreover, the law operates at two hierarchical levels:
local tissue branches (Glagov, Wolff, Hebb, Davis,
and thirteen others)
adapt their own impedance, while central controllers
(cardiac pump, regulatory brain) adapt the
$Z_{\text{source}}$ or $Z_{\text{setpoint}}$ of
entire network trees.  The controller level explains
why cardiac decompensation and chronic stress cause
\emph{multi-organ} disease: the controller's
$\Gamma$ error propagates downstream to every branch.

The existence of multiple independent molecular
implementations of the same mathematical operation
is powerful evidence that the operation is selected
for by physics, not by any specific biochemical pathway.
If C2 were an accident of NMDA-receptor chemistry,
it would not appear independently in osteocytes,
endothelial cells, keratinocytes, B cells, cardiomyocytes,
baroreceptor neurons, hepatocytes, nephrons, tenocytes,
adipocytes, erythroid progenitors, and granulosa cells.

\subsection{Implications for medicine}

If tissue remodeling is C2, then disease is
$|\Gamma| \gg 0$ that the tissue cannot correct---either
because the stimulus is too fast (learning rate $\eta$
is too small: Glagov's 40\% threshold), the stimulus
has stopped (astronaut bone loss), or the molecular
machinery is broken (immunodeficiency = broken C2 in
the immune network).

This reframes chronic disease as \emph{impedance
mismatch that has exceeded the local C2 learning
capacity}.  Diabetes is $\Gamma_{\text{endocrine}}$
trapped at a nonzero value because receptor downregulation
has saturated.  Atherosclerosis is $\Gamma_{\text{vasc}}$
where Glagov remodeling has been outpaced by plaque growth.
Osteoporosis is C2 operating with insufficient
$x_{\text{pre}}$ (too little mechanical loading).

\subsection{Sex as a boundary condition on the
$\Gamma$-profile}

The fractal geometry of the dual-network skeleton
($\beta = 2^{1/3}$, $\Omega \approx 1.25$) depends only
on $(n, d)$ and is therefore identical in males and females.
Sex differences arise not in the architecture but in the
\emph{boundary conditions}: the hypothalamic--pituitary--gonadal
(HPG) axis, a sub-controller within the regulatory brain
(Section~\ref{sec:regulatory-brain}), sets sex-specific
$Z_{\text{setpoint}}$ values for multiple leaf networks.

Concretely, gonadal hormones shift the baseline impedance
and learning rate of downstream tissues:
\begin{itemize}
  \item \textbf{Muscle:} testosterone raises
    $\eta_{\text{muscle}}$, explaining higher male
    hypertrophic response at matched loading.
  \item \textbf{Adipose:} oestrogen lowers
    $Z_{\text{setpoint}}^{\text{fat}}$, favouring
    peripheral fat deposition in females.
  \item \textbf{Bone:} oestrogen withdrawal at menopause
    produces a step change in $Z_{\text{setpoint}}^{\text{bone}}$,
    abruptly increasing $\Gamma_{\text{bone}}$---the
    $\Gamma$-Net translation of postmenopausal osteoporosis.
  \item \textbf{Immune:} oestrogen enhances immune C2
    ($\eta_{\text{immune}}^{(f)} > \eta_{\text{immune}}^{(m)}$),
    which improves pathogen clearance but also raises the
    risk of over-matching (autoimmunity; female:male ratio
    $\approx 4$:$1$ in SLE).
  \item \textbf{Vascular:} pre-menopausal oestrogen
    maintains lower $\Gamma_{\text{vasc}}$ via enhanced
    endothelial C2, delaying atherosclerosis by
    $\sim$10~years relative to males.
\end{itemize}
In all cases the physics (C1--C3), the trunk topology
($\beta$, $\Omega$), and the C2 update rule are
unchanged.  Sex is a \emph{parametric shift} in the
initial and boundary conditions---specifically, in the
$Z_{\text{setpoint}}$ vector broadcast by the HPG
sub-controller---not a different dynamical law.

\subsection{Implications for engineering}

The biological C2 networks offer a design template for
self-healing engineered systems.  Any network of impedance
interfaces that implements C2 will spontaneously
minimize reflected energy, matching its impedances to
changing loads without external tuning.  This principle
applies to self-reconfiguring antenna arrays, adaptive
power grids, and robotic systems with compliant joints.

\subsection{What is not claimed}

\begin{enumerate}
  \item We do not claim that C2 is the only molecular
    process in tissue remodeling.  Tissues undergo
    many processes (apoptosis, angiogenesis, fibrosis)
    that are not directly C2.  We claim only that
    the \emph{net effect} of remodeling on the
    tissue's impedance follows C2 on average, because
    evolution selects for energy-efficient organisms.

  \item We do not claim exact quantitative equivalence.
    Each tissue has different $\eta$ (months for bone,
    milliseconds for synapses), and the relationship
    between $Z$ and the physical variables may be
    nonlinear.  C2 is a first-order approximation valid
    in the small-$\Gamma$ regime; large perturbations
    may require higher-order terms.

  \item We do not claim that the NHANES AUC values
    prove C2 causation.  They demonstrate a prediction
    consistent with network propagation.  Causal proof
    would require interventional studies (e.g., lowering
    SBP and measuring cardiac $\Gamma$ response).
\end{enumerate}


% ============================================================
%  SECTION 8 — CONCLUSION
% ============================================================
\section{Conclusion}
\label{sec:conclusion}

We have established three results.

\textbf{First}, constraint C2
($\Delta Z = -\eta\,\Gamma\,x_{\text{pre}}\,x_{\text{post}}$)
is not a postulate but the unique gradient-descent rule
for minimizing reflected energy at a repeatedly stimulated
interface.

\textbf{Second}, nineteen known remodeling and adaptation
laws---Hebb (neural), Wolff (bone), Glagov (vascular),
Davis (muscle), immune memory, keratinization (skin),
altitude adaptation (lung), intestinal adaptation (GI),
receptor regulation (endocrine), lymph node hypertrophy
(lymphatic), tubuloglomerular feedback (renal), liver
regeneration (hepatic), tendon remodeling (connective),
leptin signaling (adipose), erythropoiesis (hematopoietic),
follicular selection (reproductive), Laplace--Starling
(cardiac pump), and baroreceptor resetting (regulatory
brain)---plus the C2 null case of cartilage
($\eta \approx 0$, predicting osteoarthritis)---are
each algebraically equivalent to C2 under the appropriate
substitution of effort, flow, and impedance variables.

\textbf{Third}, NHANES validation ($n = 49{,}774$,
zero fitted parameters) confirms the network-propagation
prediction: a vascular measurement (SBP) improves
non-vascular organ predictions (cardiac $+0.166$,
neuro $+0.061$, renal $+0.024$) through inter-network
$\Gamma$-coupling, while the control organ (endocrine)
is unchanged ($\Delta = 0.000$).

\textbf{Fourth}, the two central controllers---cardiac
pump (Laplace--Starling remodeling) and regulatory brain
(baroreceptor resetting, allostatic adaptation)---also
obey C2, but at a \emph{hierarchically higher level}:
they adapt not a local branch impedance but the source
impedance $Z_{\text{source}}$ or regulatory setpoint
$Z_{\text{setpoint}}$ for entire network trees.
Controller failure thus propagates to all downstream
networks simultaneously, explaining why chronic stress
and heart failure are multi-organ diseases.

\textbf{Fifth}, cartilage serves as a falsifiable null
case: with $\eta \approx 0$ (avascular, post-mitotic
chondrocytes), C2 predicts irreversible degradation
under chronic $\Gamma > 0$.  Osteoarthritis confirms
this prediction exactly, and tissue-engineering strategies
that increase $\eta$ (chondrocyte proliferation) should
restore C2 remodeling and halt progression.

These results extend the $\Gamma$-Net framework from
a law of interfaces (Paper~VI) to a law of
\emph{adaptive interfaces at every hierarchical level}:
every tissue in the body not only obeys
$\Gamma = (Z_2 - Z_1)/(Z_2 + Z_1)$
but actively minimizes $\Gamma$ through the same
gradient-descent rule, from connective tissue to
central controllers, from follicular selection to
liver regeneration.

The question is no longer ``Does C2 apply to tissue~X?''
The question---given that C2 is the unique energy-minimizing
update at any interface---is:
``What stops C2 from applying to tissue~X?''

If nothing stops it, then tissue~X remodels via C2.
And nothing stops it.


% ============================================================
\begin{acknowledgments}
This work used no external funding.
NHANES data were obtained from the U.S.\ Centers for Disease
Control and Prevention (CDC) National Center for Health
Statistics.
The author thanks Claude (Anthropic) for extended discussions
that crystallized the universality of C2 across tissue types.
\end{acknowledgments}


% ============================================================
\appendix

\section{Experimental Code and Data}
\label{app:code}

\begin{itemize}
  \item Network propagation experiment:
    \texttt{experiments/exp\_expanded\_physics.py}
  \item Results: \texttt{nhanes\_results/expanded\_physics\_results.json}
  \item Full calibration pipeline:
    \texttt{experiments/exp\_full\_calibration.py}
  \item Stepwise organ selection:
    \texttt{experiments/exp\_stepwise\_calibration.py}
\end{itemize}

Source code: \url{https://github.com/cyhuang76/alice-gamma-net}

Concept DOI:
\href{https://doi.org/10.5281/zenodo.18751831}{10.5281/zenodo.18751831}

Code DOI:
\href{https://doi.org/10.5281/zenodo.18852835}{10.5281/zenodo.18852835}


% ============================================================
%  BIBLIOGRAPHY
% ============================================================
\begin{thebibliography}{99}

\bibitem{PaperI}
H.-Y.~Huang, ``Irreducible Dimensional Cost in Heterogeneous
Impedance Networks,'' arXiv (2026), Paper~I of this series.

\bibitem{PaperII}
H.-Y.~Huang, ``Dual Neural--Vascular Networks from Impedance
Matching: Murray's Law, Organ Health, and the Vascular
Reflection Coefficient,'' arXiv (2026), Paper~II.

\bibitem{PaperIII}
H.-Y.~Huang, ``The Cost of Life: Impedance Debt as the
Thermodynamic Origin of Sleep and Brain Evolution,''
arXiv (2026), Paper~III.

\bibitem{PaperIV}
H.-Y.~Huang, ``The Lifecycle Equation: Cognition, Emotion,
and Aging as Three-Force Competition on an Impedance Network,''
arXiv (2026), Paper~IV.

\bibitem{PaperV}
H.-Y.~Huang, ``Memory, Consciousness, and Soul as Impedance
Physics on the $\Gamma$-Net,'' arXiv (2026), Paper~V.

\bibitem{PaperVI}
H.-Y.~Huang, ``The Grand Unification: $\Gamma$ as the
Universal Interface Law from Microtubules to Power Grids,''
arXiv (2026), Paper~VI.

\bibitem{Wolff1892}
J.~Wolff, \emph{Das Gesetz der Transformation der Knochen}
(Hirschwald, Berlin, 1892).

\bibitem{Glagov1987}
S.~Glagov, E.~Weisenberg, C.~K.~Zarins, R.~Stankunavicius,
and G.~J.~Kolettis, ``Compensatory enlargement of human
atherosclerotic coronary arteries,''
\emph{New England J.\ Med.}\ \textbf{316}, 1371 (1987).

\bibitem{Hebb1949}
D.~O.~Hebb, \emph{The Organization of Behavior}
(Wiley, New York, 1949).

\bibitem{Thomas2010}
S.~E.~Thomas, K.~Dykes, and B.~A.~Marks,
``Plantar hyperkeratosis: a study of callosities
and corns,'' \emph{J.\ Am.\ Podiatric Med.\ Assoc.}\
\textbf{100}, 251 (2010).

\bibitem{Ahmed1996}
R.~Ahmed and D.~Gray, ``Immunological memory and protective
immunity: understanding their relation,''
\emph{Science} \textbf{272}, 54 (1996).

\bibitem{Goldberg1975}
A.~L.~Goldberg, J.~D.~Etlinger, D.~F.~Goldspink, and
C.~Jablecki, ``Mechanism of work-induced hypertrophy
of skeletal muscle,''
\emph{Med.\ Sci.\ Sports} \textbf{7}, 185 (1975).

\bibitem{Frost2003}
H.~M.~Frost, ``Bone's mechanostat: a 2003 update,''
\emph{Anat.\ Record} \textbf{275A}, 1081 (2003).

\bibitem{Lang2004}
T.~Lang \emph{et al.}, ``Cortical and trabecular bone
mineral loss from the spine and hip in long-duration
spaceflight,'' \emph{J.\ Bone Miner.\ Res.}\ \textbf{19},
1006 (2004).

\bibitem{Victora2012}
G.~D.~Victora and M.~C.~Nussenzweig, ``Germinal centers,''
\emph{Annu.\ Rev.\ Immunol.}\ \textbf{30}, 429 (2012).

\bibitem{Cepeda2006}
N.~J.~Cepeda \emph{et al.}, ``Distributed practice in
verbal recall tasks: a review and quantitative synthesis,''
\emph{Psychol.\ Bull.}\ \textbf{132}, 354 (2006).

\bibitem{Hackett2001}
P.~H.~Hackett and R.~C.~Roach, ``High-altitude illness,''
\emph{New England J.\ Med.}\ \textbf{345}, 107 (2001).

\bibitem{Tappenden2014}
K.~A.~Tappenden, ``Intestinal adaptation following resection,''
\emph{J.\ Parenter.\ Ent.\ Nutr.}\ \textbf{38}, 23S (2014).

\bibitem{Kahn2006}
S.~E.~Kahn, R.~L.~Hull, and K.~M.~Utzschneider,
``Mechanisms linking obesity to insulin resistance
and type~2 diabetes,''
\emph{Nature} \textbf{444}, 840 (2006).

\bibitem{Davis1867}
H.~E.~Davis,
``Conservative Surgery,''
(Appleton, New York, 1867).

\bibitem{Hill1938}
A.~V.~Hill,
``The heat of shortening and the dynamic constants
of muscle,''
\emph{Proc.\ R.\ Soc.\ Lond.\ B} \textbf{126}, 136 (1938).

\bibitem{Schoenfeld2017}
B.~J.~Schoenfeld, J.~Grgic, D.~Ogborn, and
J.~W.~Krieger,
``Strength and hypertrophy adaptations between
low- vs.\ high-load resistance training: a systematic
review and meta-analysis,''
\emph{J.\ Strength Cond.\ Res.}\ \textbf{31}, 3508 (2017).

\bibitem{West2012}
J.~B.~West,
``High-altitude medicine,''
\emph{Am.\ J.\ Respir.\ Crit.\ Care Med.}\
\textbf{186}, 1229 (2012).

\bibitem{Brenner1982}
B.~M.~Brenner, T.~W.~Meyer, and T.~H.~Hostetter,
``Dietary protein intake and the progressive nature of
kidney disease: the role of hemodynamically mediated
glomerular injury in the pathogenesis of progressive
glomerular sclerosis in aging, renal ablation, and
intrinsic renal disease,''
\emph{New England J.\ Med.}\ \textbf{307}, 652 (1982).

\bibitem{Michalopoulos2007}
G.~K.~Michalopoulos,
``Liver regeneration,''
\emph{J.\ Cell.\ Physiol.}\ \textbf{213}, 286 (2007).

\bibitem{Friedman2019}
J.~M.~Friedman,
``Leptin and the endocrine control of energy balance,''
\emph{Nature Metab.}\ \textbf{1}, 754 (2019).

\bibitem{Buckwalter2004}
J.~A.~Buckwalter, H.~J.~Mankin, and A.~J.~Grodzinsky,
``Articular cartilage and osteoarthritis,''
\emph{Instr.\ Course Lect.}\ \textbf{54}, 465 (2005).

\bibitem{Grossman1975}
W.~Grossman, D.~Jones, and L.~P.~McLaurin,
``Wall stress and patterns of hypertrophy in the human
left ventricle,''
\emph{J.\ Clin.\ Invest.}\ \textbf{56}, 56 (1975).

\bibitem{Maron2006}
B.~J.~Maron and A.~Pelliccia,
``The heart of trained athletes: cardiac remodeling and the
risks of sports, including sudden death,''
\emph{Circulation} \textbf{114}, 1633 (2006).

\bibitem{Chapleau1991}
M.~W.~Chapleau, G.~Hajduczok, and F.~M.~Abboud,
``Mechanisms of resetting of arterial baroreceptors:
an overview,''
\emph{Am.\ J.\ Med.\ Sci.}\ \textbf{301}, 271 (1991).

\bibitem{McEwen1998}
B.~S.~McEwen, ``Stress, adaptation, and disease:
Allostasis and allostatic load,''
\emph{Ann.\ N.Y.\ Acad.\ Sci.}\ \textbf{840}, 33 (1998).

\end{thebibliography}

\end{document}
